이 책은 사람에게는 당연해 보이는 계산 하나를 컴퓨터로 수행하려면 무엇을 결정해야 하며, 그 결정마다 어떤 비용이 따르는지를 묻는다.
같은 답을 얻는 계산도 작은 계산으로 어떻게 나눌지, 어떤 순서로 실행할지, 중간값을 어디에 둘지에 따라 걸리는 시간과 필요한 회로, 데이터 이동량이 달라진다.

우리는 진리표에서 출발해 작은 논리 회로와 덧셈기·곱셈기를 만들고, 하나의 CPU 코어가 계산을 수행하는 원리를 살펴본다.
이어 여러 계산을 동시에 처리하는 SIMD와 SIMT, 워프(warp), GPU 구조로 올라가면서 계산을 복제하고 자원을 공유하는 방법을 배운다.

행렬 곱셈은 처음부터 되풀이하는 예제가 아니라 이 여정의 종합 문제로 사용한다.
앞에서 배운 계산의 분해, 실행 순서, 병렬 처리, 데이터 재사용이 하나의 실제 계산 안에서 어떻게 만나며, 그 과정에서 어떤 비용이 발생하는지를 보여줄 것이다.
